\documentclass[book.tex]{subfiles}

\begin{document}
\chapter{Optimization}
\label{chap:optimization}
\noindent\rule{11cm}{0.4pt} \\
\textbf{Prerequisites:} \autoref{chap:calculus} \\
\textbf{Difficulty Level:} ** \\
\noindent\rule{11cm}{0.4pt}
\vspace{5mm}

Both machine learning and physics require methods to find points of functions that satisfy certain maximal or minimal criteria: for example the "highest" point on a 1D function. Optimization is the field of mathematics that provides procedures to find such special points. In the case that the functions in question take values on function spaces, then optimization turns into the problem of finding a function with a desired property, a technique of powerful physical importance.

\section{Optimizing Simple Functions}

As a simple example, let us suppose that we have a function $f: \mathbb{R} \to \mathbb{R}$. How can we find the extrema of $f$? The first place to start looking is at the zeros of the derivative $f'$. For example, suppose that $f(x) = x^2$, then $f'(x) = 2x$ which has a zero at $x=0$. This is the minima of $f$. This rule isn't perfect though. Suppose $g(x) = x^3$. Then $g'(x) = 3x^2$ which also has a zero at $x=0$, but this isn't a minima of $g$. 


\section{Lagrange Multipliers}

Lagrange multipliers are a technique for finding minima that satisfy constraints. For example, suppose that $f, g: \mathbb{R}^3 \to \mathbb{R}$ are functions of three variables. Then we would like to solve the problem.
\begin{align}
    \min &\quad f(x, y, z)\\
    \textrm{s. t.} &\quad g(x, y, z) = k
\end{align}

The technique for doing this is to find a spot such that the gradients of $f$ and $g$ point in the same direction and where the constraint is satisified. That is, we construct the following set of equations.

\begin{align}
    \nabla f(x, y, z) &= \nabla g(x, y, z) \\
    g(x, y, z) &= k
\end{align}

Then check all solutions of this set of equations for those which are minima/maxima of the original equation. As a singularity check, also verify that the gradient of $g$ is not zero at that point either.

\section{Adjoint Method}

Let $u \in \mathcal{U}$ be a state variable. Let $v \in \mathcal{V}$ be an optimization variable. Let
\begin{align}
J: \mathcal{U} \times \mathcal{V} \to \mathbb{R}
\end{align}
be a function and considered the objective function
\begin{align}
    j(v) &= J(u(v), v)
\end{align}
We can then formulate the optimization problem
\begin{align}
    & \textrm{minimize}\ j(v) = J(u_v, v) \\
    & \textrm{subject to}\ D_v(u) = 0
\end{align}
We wish to compute the total derivative $d_\theta j$ subject to the constraint $D_v(u) = 0$. Create the Lagrangian
\begin{align}
    \mathcal{L}(u, v, \lambda) &= J(u, v) + \lambda D_v(u) \rangle 
\end{align}
The method of Lagrange multipliers means solution has to be a stationary point of Lagrangian. The adjoint then comes out
\begin{align}
d_\theta f(x) &= \lambda^T g_\theta
\end{align}

\section{Convex Optimization}

Convex problems can be exactly solved by gradient descent and related optimization methods.

\section{Calculus of Variations}

Differential calculus is concerned with the task of finding the extreme values of functions of one or more variables. This can be done by finding points where the derivative of the function is 0. However, in some cases, we will be interested in finding an extreme point in a set of functions. 

The \textit{brachistochrone problem} is a famous example of such a task. This challenge, first posed in 1692, asks for the curve between two points $A$, $B$ through which a particle acted upon only by gravity can move from $A$ to $B$. In order for us to be able to solve this problem, we will have to introduce some machinery. 

Let $\mathcal{F}$ be a space of functions. A function $J: \mathcal{F} \to \mathbb{R}$ is said to be a \textit{functional}. We denote the evaluation of $J$ on $f \in \mathcal{F}$ as $J[f]$.  We say that $J$ has an local maximum at $f$ if for all $y$ which are in a neighborhood of $f$,  $J[y] \geq J[f]$. Alternatively, if $J[y] \leq J[f]$ for all such $y$, we say $J$ has a local minimum

If you haven't worked with functionals before, this definition might be a little opaque. Let's look at an example. Let $x_1, x_2 \in \mathbb{R}$ be constants. Let $y: \mathbb{R} \to \mathbb{R}$ be a twice differentiable function. Let $L: \mathbb{R}^3 \to \mathbb{R}$ be a twice differentiable function. Then we define
\begin{equation}
    J[f] = \int_{x_1}^{x_2} L(x, f(x), f'(x)) dx
\end{equation}

Here $f' = \frac{df}{dx}$. How can we find the minima of $J$? We typically do so by solving the \textit{Euler-Lagrange equations}.
\begin{equation}
    \frac{\partial L}{\partial f} = \frac{d}{dx} \frac{\partial L}{\partial f'}
\end{equation}

In the exercises, we'll provide some examples of using the Euler-Lagrange equations on practice problems

%\textcolor{red}{TODO: Define the functional derivative}

\subsection{Exercises}
\begin{enumerate}
    \item Let $f(x) = 4x^2 - 3y$. Find the extrema of $f$ on the circle $x^2 + y^2 = 4$.
    \item Let $(x_1, y_1), (x_2, y_2) \in \mathbb{R}^2$ be tuples of points. We would like to to find the shortest path that connects $(x_1, y_1)$ and $(x_2, y_2)$. We define a functional that computes the arc-length of a curve.
    \begin{equation}
        A[f] = \int_{x_1}^{x_2} \sqrt{1 + (f'(x))^2} dx
    \end{equation}
    \begin{enumerate}
        \item Let $g$ be the straight line that connects $(x_1, y_1)$ and $(x_2, y_2)$. Write down the equation of $g$.
        \item Compute $A[g]$. Try a few other paths between these points and convince yourself that the arclength is smallest for $g$. 
        \item Show that $\frac{\partial L}{\partial f} = 0$. Hint: Note that $f$ and $f'$ are independent arguments for $L$.
        \item Show that 
        \begin{equation}
            \frac{\partial L}{\partial f'} = \frac{f'(x)}{\sqrt{1 + (f'(x))^2}}
        \end{equation}
        \item Show that 
        \begin{equation}
            \frac{f'(x)}{\sqrt{1 + (f'(x))^2}} = c
        \end{equation}
        where $c \in \mathbb{R}$ is a constant. Hint: plug the results of the previous sections into the Euler-Lagrange equations.
        \item Show that $f'(x)$ must be a constant. Prove that this implies $f = g$
    \end{enumerate}
\end{enumerate}


\end{document}
